% Options for packages loaded elsewhere
\PassOptionsToPackage{unicode}{hyperref}
\PassOptionsToPackage{hyphens}{url}
\documentclass[
  12pt,
]{ctexart}
\usepackage{xcolor}
\usepackage{amsmath,amssymb}
\setcounter{secnumdepth}{-\maxdimen} % remove section numbering
\usepackage{iftex}
\ifPDFTeX
  \usepackage[T1]{fontenc}
  \usepackage[utf8]{inputenc}
  \usepackage{textcomp} % provide euro and other symbols
\else % if luatex or xetex
  \usepackage{unicode-math} % this also loads fontspec
  \defaultfontfeatures{Scale=MatchLowercase}
  \defaultfontfeatures[\rmfamily]{Ligatures=TeX,Scale=1}
\fi
\usepackage{lmodern}
\ifPDFTeX\else
  % xetex/luatex font selection
\fi
% Use upquote if available, for straight quotes in verbatim environments
\IfFileExists{upquote.sty}{\usepackage{upquote}}{}
\IfFileExists{microtype.sty}{% use microtype if available
  \usepackage[]{microtype}
  \UseMicrotypeSet[protrusion]{basicmath} % disable protrusion for tt fonts
}{}
\makeatletter
\@ifundefined{KOMAClassName}{% if non-KOMA class
  \IfFileExists{parskip.sty}{%
    \usepackage{parskip}
  }{% else
    \setlength{\parindent}{0pt}
    \setlength{\parskip}{6pt plus 2pt minus 1pt}}
}{% if KOMA class
  \KOMAoptions{parskip=half}}
\makeatother
\usepackage{longtable,booktabs,array}
\usepackage{calc} % for calculating minipage widths
% Correct order of tables after \paragraph or \subparagraph
\usepackage{etoolbox}
\makeatletter
\patchcmd\longtable{\par}{\if@noskipsec\mbox{}\fi\par}{}{}
\makeatother
% Allow footnotes in longtable head/foot
\IfFileExists{footnotehyper.sty}{\usepackage{footnotehyper}}{\usepackage{footnote}}
\makesavenoteenv{longtable}
\setlength{\emergencystretch}{3em} % prevent overfull lines
\providecommand{\tightlist}{%
  \setlength{\itemsep}{0pt}\setlength{\parskip}{0pt}}
\usepackage{bookmark}
\IfFileExists{xurl.sty}{\usepackage{xurl}}{} % add URL line breaks if available
\urlstyle{same}
\hypersetup{
  hidelinks,
  pdfcreator={LaTeX via pandoc}}

\author{SCX}
\date{2026}

\begin{document}

\section{Tomorrow's TODOs TODO ---
2026-06-29}\label{ux660eux65e5ux63a8ux8fdb-todo-2026-06-29}

\subsection{🔴 Full Drug Database下载 + Yajie
One-Shot Full Screening}\label{ux5168ux836fux7269ux6570ux636eux5e93ux4e0bux8f7d-yajie-ux4e00ux6b21ux6027ux5168ux7b5b}

\subsubsection{Database List (Full, not limited to
HIV）}\label{ux6570ux636eux5e93ux6e05ux5355ux5168ux91cfux4e0dux4ec5ux9650ux4e8e-hiv}

\begin{longtable}[]{@{}
  >{\centering\arraybackslash}p{(\linewidth - 8\tabcolsep) * \real{0.1667}}
  >{\raggedright\arraybackslash}p{(\linewidth - 8\tabcolsep) * \real{0.2500}}
  >{\raggedright\arraybackslash}p{(\linewidth - 8\tabcolsep) * \real{0.2500}}
  >{\centering\arraybackslash}p{(\linewidth - 8\tabcolsep) * \real{0.1667}}
  >{\centering\arraybackslash}p{(\linewidth - 8\tabcolsep) * \real{0.1667}}@{}}
\toprule\noalign{}
\begin{minipage}[b]{\linewidth}\centering
\#
\end{minipage} & \begin{minipage}[b]{\linewidth}\raggedright
数据库
\end{minipage} & \begin{minipage}[b]{\linewidth}\raggedright
内容
\end{minipage} & \begin{minipage}[b]{\linewidth}\centering
大小
\end{minipage} & \begin{minipage}[b]{\linewidth}\centering
优先级
\end{minipage} \\
\midrule\noalign{}
\endhead
\bottomrule\noalign{}
\endlastfoot
1 & \textbf{ChEMBL} & 全量生物活性数据（所有靶点、所有化合物） &
\textasciitilde50GB & 🔴 \\
2 & \textbf{DrugBank} & 全部已批准+实验性药物，带靶点 &
\textasciitilde1GB & 🔴 \\
3 & \textbf{PubChem BioAssay} & 全部检测数据 & \textasciitilde50GB &
🔴 \\
4 & \textbf{BindingDB} & 全量蛋白-配体结合亲和力 & \textasciitilde2GB &
🔴 \\
5 & \textbf{PDBbind} & 蛋白-配体复合物+结合亲和力（最精确） &
\textasciitilde5GB & 🟡 \\
6 & \textbf{TTD} (Therapeutic Target Database) & 已知治疗靶点+对应药物 &
\textasciitilde500MB & 🔴 \\
7 & \textbf{DrugCentral} & FDA批准药物的靶点+适应症 & \textasciitilde1GB
& 🟡 \\
8 & \textbf{Open Targets} & 靶点-疾病关联（GWAS+文献） &
\textasciitilde10GB & 🟡 \\
9 & \textbf{PharmGKB} & 药物基因组学（基因×药物×反应） &
\textasciitilde2GB & 🟡 \\
10 & \textbf{Stanford HIVDB} & HIV耐药突变（保留） & \textless100MB &
🔴 \\
11 & \textbf{SIDER} & 药物副作用数据库 & \textasciitilde1GB & 🟡 \\
12 & \textbf{STITCH} & 化合物-蛋白相互作用网络 & \textasciitilde20GB &
🟢 \\
\end{longtable}

\subsubsection{硬盘需求}\label{ux786cux76d8ux9700ux6c42}

\begin{verbatim}
🔴 优先级 (必须): ~105GB
🟡 强推荐:         ~20GB  
🟢 可选:           ~20GB
缓冲+解压:        ~50GB
─────────────────────────
总计:             ~200GB
\end{verbatim}

\subsubsection{预处理管道}\label{ux9884ux5904ux7406ux7ba1ux9053}

\begin{itemize}
\tightlist
\item[$\boxtimes$]
  所有 SMILES → 标准化 → ECFP4/MACCS 指纹 （screen\_all\_databases.py
  Phase 2）
\item[$\boxtimes$]
  统一 schema: drug\_id, target\_id, assay\_type, value, units,
  source\_db
\item[$\boxtimes$]
  去重（同一 drug-target 对在多个数据库中重复出现 = 专家共识）
\end{itemize}

\subsubsection{Yajie 运行}\label{yajie-ux8fd0ux884c}

\begin{itemize}
\tightlist
\item[$\boxtimes$]
  N 个专家（每数据库一个 + 物理化学性质 + 结构相似性）
\item[$\boxtimes$]
  一次性筛选：所有药物 × 所有靶点 （screen\_all\_databases.py）
\item[$\boxtimes$]
  输出：

  \begin{itemize}
  \tightlist
  \item
    高共识 drug-target 对（多数据库交叉验证） → mt\_high\_confidence.csv
  \item
    数据库质量报告（哪个数据库矛盾最多） → mt\_quality\_report.json
  \item
    新兴靶点（高新颖性低共识 → Spring 待复活） →
    mt\_novel\_candidates.csv
  \end{itemize}
\end{itemize}

\subsubsection{验证}\label{ux9a8cux8bc1}

\begin{itemize}
\tightlist
\item[$\square$]
  已知 HIV 药物是否被 Yajie 高分筛出
\item[$\square$]
  Yajie 是否发现已知的跨界靶点（HIV/自身免疫/炎症共靶点）
\item[$\square$]
  False Positive分析（高分但文献无支持的 → 新发现候选）
\end{itemize}

\begin{center}\rule{0.5\linewidth}{0.5pt}\end{center}

\subsection{🟡 arXiv 投稿}\label{arxiv-ux6295ux7a3f}

\begin{itemize}
\tightlist
\item[$\square$]
  Paper 1 (Yajie) --- PDF已编译（含S1-S8）
\item[$\square$]
  Paper 2 (Spring) --- PDF已编译（13页）
\end{itemize}

\begin{center}\rule{0.5\linewidth}{0.5pt}\end{center}

\subsection{✅ Completed}\label{ux5df2ux5b8cux6210}

\begin{itemize}
\tightlist
\item[$\boxtimes$]
  Spring 代码 (1551行) + Yajie fit() + HIV 审计原型
\item[$\boxtimes$]
  9篇论文框架 + 18层博弈论 + 445测试全过
\item[$\boxtimes$]
  \textbf{Full Drug Database下载脚本}
  \texttt{drug-module/scripts/download\_databases.py}
  (\textasciitilde1050行)

  \begin{itemize}
  \tightlist
  \item
    12个数据库完整覆盖（ChEMBL, DrugBank, PubChem BioAssay, BindingDB,
    TTD, Stanford HIVDB, PDBbind, DrugCentral, Open Targets, PharmGKB,
    SIDER, STITCH）
  \item
    分层下载（Tier 1/2/3）约205GB
  \item
    每个数据库写入 provenance.json（SHA256校验、下载日期、Version、URL）
  \item
    进度条+ETA、断点续传、认证数据库占位符
  \item
    自动化解析: ChEMBL SQLite→Parquet, DrugBank XML→Parquet, BindingDB
    TSV→Parquet, PubChem PUG REST API, Open Targets GraphQL API
  \end{itemize}
\item[$\boxtimes$]
  \textbf{Yajie全数据库筛查脚本}
  \texttt{drug-module/scripts/screen\_all\_databases.py}
  (\textasciitilde930行)

  \begin{itemize}
  \tightlist
  \item
    Phase 1: 加载所有数据库 Parquet/CSV + provenance 元数据
  \item
    Phase 2: SMILES→ECFP4指纹标准化 + 靶点→UniProt ID映射
  \item
    Phase 3: 药物-靶点统一证据矩阵（行=药靶对, 列=数据库专家）
  \item
    Phase 4: Yajie多专家共识评分（Theorem 1: CLEAN/NOISY/AMBIGUOUS分类）
  \item
    Phase 5: 8个输出files（MT Gold Standard CSV, 分类Abstract,
    数据库Consistency矩阵, 高置信度对, Disagreement报告, 新候选, 质量报告JSON,
    溯源审计JSON）
  \item
    完全可溯源：每对记录包含source\_databases, database\_links,
    download\_date, mt\_report\_version
  \end{itemize}
\end{itemize}

\end{document}
